% ============================================================
% 山东省大学生数学建模竞赛论文模板
% ============================================================
% 使用说明：
% 1. 填写下方的团队信息
% 2. 在 chapters/ 目录下编辑各章节内容
% 3. 使用 tectonic -X compile main.tex 编译
% ============================================================

% 文档类选项：
% - withoutpreface: 去掉封面与编号页（电子版提交时使用）
% - bwprint: 黑白打印模式
\documentclass[withoutpreface,bwprint]{main}

% 额外宏包
\usepackage[framemethod=TikZ]{mdframed}
\usepackage{url}
\usepackage{subcaption}

% ============================================================
% 论文基本信息设置
% ============================================================
\title{山东省大学生数学建模竞赛论文模板}

% 团队信息（取消注释并填写）
% \tihao{A}                      % 题号（A/B/C/D/E）
% \baominghao{xxxx}              % 报名号（12位数字）
% \schoolname{XX大学}            % 学校全称
% \membera{成员A}                % 队员1
% \memberb{成员B}                % 队员2
% \memberc{成员C}                % 队员3
% \supervisor{指导教师}          % 指导教师
% \yearinput{2026}               % 年
% \monthinput{09}                % 月
% \dayinput{04}                  % 日

% 提取论文标题，用于摘要页输出
\makeatletter
\let\thesistitle\@title
\makeatother

\begin{document}

% ============================================================
% 第1页：目录（无页码）
% ============================================================
\thispagestyle{empty}
\tableofcontents
\newpage

% ============================================================
% 第2页：摘要专用页（页码从1开始）
% ============================================================
\setcounter{page}{1}
\begin{center}
{\zihao{3}\bfseries \thesistitle\par}
\end{center}
\vskip1ex
% ============================================================
% 摘要
% ============================================================

\begin{abstract}
本模板是为山东省大学生数学建模竞赛编写的 \LaTeX{} 模板，旨在让大家专注于论文的内容写作，而不用花费过多精力在格式的定制和调整上。本手册是相应的参考，其中提供了一些环境和命令可以让模板的使用更为方便。同时需要注意，使用者需要有一定的 \LaTeX{} 的使用经验，至少要会使用常用宏包的一些功能，比如参考文献、数学公式、图片使用、列表环境等等。

\keywords{\TeX{}\quad 图片\quad 表格\quad 公式}
\end{abstract}


% ============================================================
% 第3页开始：正文
% ============================================================
% ============================================================
% 第一章 问题重述
% ============================================================

\section{问题重述}

\subsection{问题背景}

此处简要描述问题的实际背景，说明问题来源于哪个领域，具有什么样的实际意义与研究价值。

\subsection{问题的提出}

根据竞赛题目，将原问题重述如下：

\begin{problem}
	问题的具体描述。可以分小点列出各问。
	\label{pro:example}
\end{problem}

针对上述问题，本文将依次展开分析与建模。
      % 第一章 问题重述
% ============================================================
% 第二章 问题分析
% ============================================================

\section{问题分析}

\subsection{问题一的分析}

针对问题一，我们首先对题目条件进行梳理，明确已知数据和待求解目标。总体思路如图~\ref{fig:analysis-1}所示。

\begin{figure}[!h]
    \centering
    \includegraphics[width=.6\textwidth]{example}
    \caption{问题一分析流程图}
    \label{fig:analysis-1}
\end{figure}

\subsection{问题二的分析}

针对问题二，我们需要在问题一的基础上进一步拓展。关键在于建立更一般的模型或引入新的约束条件。

\subsection{问题三的分析}

针对问题三，综合前两问的结果，进行敏感性分析或推广讨论。
      % 第二章 问题分析
% ============================================================
% 第三章 模型假设
% ============================================================

\section{模型假设}

为了对问题进行合理的简化和抽象，使模型更具可操作性，本文提出以下模型假设：

\begin{enumerate}
	\item 假设一的具体内容；
	\item 假设二的具体内容；
	\item 假设三的具体内容；
	\item 假设四的具体内容。
\end{enumerate}

此外，本文假设所有数据真实可靠，且在研究的时间范围内，相关环境参数保持稳定。
      % 第三章 模型假设
% ============================================================
% 第四章 符号说明
% ============================================================

\section{符号说明}

本文涉及的主要符号及其说明如\cref{tab:symbols}所示。

\begin{table}[!htbp]
	\caption{符号说明}\label{tab:symbols} \centering
	\begin{tabular}{ccc}
		\toprule[1.5pt]
		符号 & 含义 & 单位 \\
		\midrule[1pt]
		$D$ & 木条宽度 & cm \\
		$P$ & 功率 & W \\
		$E$ & 能量 & J \\
		\bottomrule[1.5pt]
	\end{tabular}
\end{table}

注：其他符号在文中使用时另行说明。
      % 第四章 符号说明
% ============================================================
% 第五章 模型的建立与求解
% ============================================================

\section{模型的建立与求解}

\subsection{问题一的模型建立与求解}

\subsubsection{模型建立}

数学建模必然涉及不少数学公式的使用。如果希望某个公式带编号，并且在后文中引用可以参考下面的写法：
\begin{equation}
	E=mc^2
	\label{eq:energy}
\end{equation}
式\cref{eq:energy}是质能方程。

多行公式有时候希望能够在特定的位置对齐，以下是其中一种处理方法。
\begin{align}
	P & = UI \\
	& = I^2R
\end{align}

矩阵的输入也不难。
\[
\mathbf{X} = \left(
    \begin{array}{cccc}
    x_{11} & x_{12} & \ldots & x_{1n}\\
    x_{21} & x_{22} & \ldots & x_{2n}\\
    \vdots & \vdots & \ddots & \vdots\\
    x_{n1} & x_{n2} & \ldots & x_{nn}\\
    \end{array} \right)
\]

分段函数这些可以用 \verb|cases| 环境，但是它要放在数学环境里面。
\[
f(x) =
    \begin{cases}
        0 &  x \text{为无理数} ,\\
        1 &  x \text{为有理数} .
    \end{cases}
\]

\begin{theorem}
	这是一个定理。
	\label{thm:example}
\end{theorem}

\begin{proof}
	这是一个证明。
\end{proof}

\subsubsection{模型求解}

描述求解算法与步骤，可配合流程图或伪代码进行说明。

\subsection{问题二的模型建立与求解}

\subsection{问题三的模型建立与求解}
      % 第五章 模型的建立与求解

% ============================================================
% 模型评价
% ============================================================
% ============================================================
% 第六章 模型的优缺点
% ============================================================

\section{模型的优缺点}

\subsection{模型优点}

\begin{enumerate}
	\item 优点一；
	\item 优点二；
	\item 优点三。
\end{enumerate}

\subsection{模型缺点}

\begin{enumerate}
	\item 缺点一；
	\item 缺点二。
\end{enumerate}

\subsection{模型改进方向}

针对上述缺点，可以从以下方面进行改进：

\begin{enumerate}
	\item 改进方向一；
	\item 改进方向二。
\end{enumerate}
    % 第六章 模型的优缺点

% ============================================================
% 参考文献
% ============================================================
% ============================================================
% 第七章 参考文献
% ============================================================

\begin{thebibliography}{9}
    \bibitem[1]{liuhaiyang2013latex}
    刘海洋.
    \newblock \LaTeX {}入门\allowbreak[J].
    \newblock 电子工业出版社, 北京, 2013.
    \bibitem[2]{mathematical-modeling}
    山东省大学生数学建模竞赛论文格式规范.
\end{thebibliography}
    % 第七章 参考文献

% ============================================================
% 附录
% ============================================================
\input{chapters/appendix}      % 第八章 附录

\end{document}
